\input{../tex-inputs/latex-leseansicht-vorspann}

\section[Gustav Schwarzkopf an Arthur Schnitzler, 1{[}0?{]}. 9. 1927]{L04242 Gustav Schwarzkopf an Arthur Schnitzler, 1{[}0?{]}. 9. 1927}
\nopagebreak\mylabel{L04242v}
\rehead{ }\normalsize\beginnumbering\briefempfaengerindex{Schnitzler, Arthur@\textsc{Schnitzler, Arthur}!zzzSchwarzkopf, Gustav@\emph{von Gustav Schwarzkopf}!1927-09-101@{1{[}0?{]}. 9. 1927}|(be}
\toendnotes[C]{\smallbreak\pagebreak[2]}
\correspDesc{Versand  durch Gustav Schwarzkopf am 1[0?]. 9. 1927 in Wien
\newline{}Übermittlung  am 12. 9. 1927 in Wien
\newline{}Erhalt  durch Arthur Schnitzler im Zeitraum [13. 9. 1927 – 15. 9. 1927?] in Venedig}\toendnotes[C]{\smallbreak}
\Standort{CUL, Schnitzler, B 96a.}
\physDesc{Postkarte, 632 Zeichen
\newline{}Handschrift: schwarze Tinte, deutsche Kurrent
\newline{}Versand: Stempel: »\nobreak{}\oindex{I., Innere Stadt@\textbf{I., Innere Stadt}, \emph{Verwaltungsgebiet}|pwk}1/\textsubscript{1} Wien 9, 12. IX. 27, 14\nobreak{}«.  }\toendnotes[C]{\smallbreak}\pstart{}{\pb}I.
                     Tiefer Graben 17. II. 6\oindex{Wien@\textbf{Wien}!I., Innere Stadt@\textbf{I., Innere Stadt}!Tiefer Graben 17@\textbf{Tiefer Graben 17}, \emph{Wohngebäude}|pw}\pend{}{\bigskip}\pstart{}Herr\pend{}\pstart{}Dr Arthur Schnitzler\pend{}\pstart{}Venedig Lido\oindex{Lido@\textbf{Lido}|pw}\pend{}\pstart{}Corno d’oro\oindex{Hotel Corno d’Oro@\textbf{Hotel Corno d’Oro}, \emph{Hotel}|pw}\pend{}{\bigskip}\vspace{1em}
\pstart
           \raggedleft{}{\pb}1\textcolor{gray}{0}. IX 27\pend
           \vspace{0.5em}
\pstart
           Lieber Arthur, ich danke Ihnen für Ihr Karten – beſonders für die
               \label{K_L04242-1v}\edtext{letzte}{\lemma{\textnormal{\emph{letzte}}}\Cendnote{\textnormal{{XXXX ref} XXXX
               }}}\label{K_L04242-1} – die Arena in Verona\oindex{Verona Arena@\textbf{Verona Arena}, \emph{Amphitheater}|pw} – da war ich auch einmal, es{ }ſind nur 32 Jahre. –
               Meinem Bruder Max\pwindex{Schwarzkopf, Max 12.\,6.\,1857 Wien – 14.\,4.\,1928 ebd.@\textsc{Schwarzkopf, Max} (12.\,6.\,1857 Wien – 14.\,4.\,1928 ebd.), \emph{Rechtsanwalt}|pw} und mir geht es gut, man{ }ſagt uns,
               daß wir{ }ſo ausſehen, als ob wir eben vom Urlaub kämen, aber auch Max\pwindex{Schwarzkopf, Max 12.\,6.\,1857 Wien – 14.\,4.\,1928 ebd.@\textsc{Schwarzkopf, Max} (12.\,6.\,1857 Wien – 14.\,4.\,1928 ebd.), \emph{Rechtsanwalt}|pw} iſt über Schönbru{\geminationn}\oindex{Wien@\textbf{Wien}!XIII., Hietzing@\textbf{XIII., Hietzing}!Schlosspark Schönbrunn@\textbf{Schlosspark Schönbrunn}, \emph{Park}|pw}
                und Baden\oindex{Baden bei Wien@\textbf{Baden bei Wien}, \emph{Hauptstadt}|pw} nicht hinaus
                  geko{\geminationm}en.\pend
           
\pstart
           – – Bitte, beſtellen Sie gütigſt Frau \label{K_L04242-2v}\edtext{Lili\pwindex{Cappellini, Lili 13.\,9.\,1909 Wien – 26.\,7.\,1928 Venedig@\textsc{Cappellini, Lili} (13.\,9.\,1909 Wien – 26.\,7.\,1928 Venedig)|pw} meine
               ergebenſten Grüße und Glückwünſche zum \uline{18}.
               Geburtstag}{\lemma{\textnormal{\emph{Lili … Geburtstag}}}\Cendnote{\textnormal{am 13. 9. 1927}}}\label{K_L04242-2}.\pend
           
\pstart
           {\pb}Sie bleiben wol noch länger? Na, da Sie{ }ſchon die \label{K_L04242-3v}\edtext{erſte Generalprobe\eventindex{Burgtheater@\textbf{Burgtheater}!Generalprobe von Im Wirtshaus zum Pechvogel, 10.9.1927@Generalprobe von Im Wirtshaus zum Pechvogel, 10.9.1927|pwv} im Burgtheater\orgindex{Burgtheater@Burgtheater|pw} verpaßt}{\lemma{\textnormal{\emph{erste … verpaßt}}}\Cendnote{\textnormal{Schnitzler
                     besaß eine Sondergenehmigung, die ihm erlaubte, alle Generalproben des \emph{Burgtheaters}\orgindex{Burgtheater@Burgtheater|pwk} zu besuchen. Das
                     erste neu einstudierte Stück der Saison 1927/1928 war \emph{Im Wirtshaus zum Pechvogel}\pwindex{\textcolor{red}{\textsuperscript{XXXX indx1}}!Im Wirtshaus zum Pechvogel. Komödie in drei Akten@\strich\emph{Im Wirtshaus zum Pechvogel. Komödie in drei Akten}|pwk},
                     die Generalprobe\eventindex{Burgtheater@\textbf{Burgtheater}!Generalprobe von Im Wirtshaus zum Pechvogel, 10.9.1927@Generalprobe von Im Wirtshaus zum Pechvogel, 10.9.1927|pwkv} fand am 10. 9. 1927 statt.}}}\label{K_L04242-3} haben, iſt nicht mehr viel zu
               verlieren.\pend
           
\pstart
           Herzlichſt{\\[\baselineskip]}Ihr{\\[\baselineskip]}\spacefill\mbox{Gustav}.\pend
           \leftskip=0em{}\selectlanguage{ngerman}\endnumbering\briefempfaengerindex{Schnitzler, Arthur@\textsc{Schnitzler, Arthur}!zzzSchwarzkopf, Gustav@\emph{von Gustav Schwarzkopf}!1927-09-101@{1{[}0?{]}. 9. 1927}|)be}\mylabel{L04242h}
\begin{anhang}
\end{anhang}\newcommand{\dateiname}{L04242}\newcommand{\titel}{Gustav Schwarzkopf an Arthur Schnitzler, 1[0?]. 9. 1927}\newcommand{\editorInnen}{Herausgegeben von Jahnke, SelmaMüller, Martin Anton}\input{../tex-inputs/latex-leseansicht-abspann}
